% sec-03-ind.tex - Section 3, IND.  A main section.  FULL stage.
% Figures 6, 7, 8 and 9 drawn; Tables 8, 9 and 10 populated from the twelve
% section files of trial-ind/final-ind/publication.
\section{IND}
\label{sec:ind}

\subsection{What was filed, and under what authority}

The Investigational New Drug application deposited July 1, 2026 supports a Phase
1, first-in-human, combined IND and Investigational Device Exemption
investigation of on-premises large language model directed eight-arm robotic
pancreaticoduodenectomy with perioperative daraxonrasib (RMC-6236) in patients
with KRAS G12-mutated pancreatic ductal adenocarcinoma \cite{phase1ind}. The
dossier states its own authority in its opening subsection: the investigation
``is filed under 21 CFR \S 312 for the investigational drug and is coupled to a
significant-risk early-feasibility device study under 21 CFR \S 812 for the
LLM-directed robotic platform; the two pathways are coordinated within a single
integrated protocol because the device conducts the resection that generates the
operative and pathology inputs to the daraxonrasib restart advisory''
\cite{phase1ind}. That sentence is doing regulatory work rather than rhetorical
work: it is the reason a single dossier can carry both a drug and a device
without either pathway being subordinate to the other.

The clinical problem the dossier addresses is stated with the same discipline.
PDAC is the third leading cause of cancer death in the United States with total
five-year overall survival below 13 percent \cite{siegel2025statistics}, and the
pancreaticoduodenectomy requires resection and reconstruction of four luminal
structures and dissection around five named vessels. The 2025 Dutch nationwide
cohort of 1000 robotic pancreaticoduodenectomies, which anchors the contemporary
human baseline, reported a conversion rate of 10.1 percent, Clavien-Dindo grade
III or higher complications of 41.3 percent, International Study Group on
Pancreatic Surgery grade B or C fistula of 24.4 percent, and 30-day mortality of
3.9 percent \cite{dutchcohort2025, bassi2017isgps}. Figure 6 shows what the
dossier was assembled from, and Table 8 lists the twelve modules it was
assembled into.

% ---------------------------------------------------------------------------
% FIGURE 6.  diagrams (python)-type, clustered infrastructure.  Four source
% clusters in two columns, one subject node with a halo, one composed cluster.
% Every cluster contributes exactly one edge from a stated boundary waypoint.
% ---------------------------------------------------------------------------
\begin{appfloat}[!tb]
\begin{appfig}[diagrams (python)-type / clustered infrastructure]
% cluster 1, regulatory
\dgnode{dgtile}{0.75}{0.55}{\glyphdoc}{21 CFR 312 adapted}{r1}
\dgnode{dgtile}{2.95}{0.55}{\glyphdoc}{ICH E6 R3 adapted}{r2}
\dgnode{dgtile}{0.75}{-1.35}{\glyphdoc}{ReGARDD template}{r3}
\dgnode{dgtile}{2.95}{-1.35}{\glyphlock}{FDA 1571 instructions}{r4}
\begin{scope}[on background layer]
\node[dgcluster,fit=(r1)(r1l)(r2)(r2l)(r3)(r3l)(r4)(r4l)] (cr) {};\end{scope}
\node[dgctitle,anchor=south west] at ([yshift=1.5mm]cr.north west) {Regulatory source};
% cluster 2, clinical
\dgnode{dgtilem}{0.75}{-3.75}{\glyphscalpel}{Phase 1 protocol}{c1}
\dgnode{dgtilem}{2.95}{-3.75}{\glyphscalpel}{Phase 2 protocol}{c2}
\dgnode{dgtilem}{0.75}{-5.65}{\glyphchart}{Dutch 1000-case cohort}{c3}
\dgnode{dgtilem}{2.95}{-5.65}{\glyphpill}{RASolute 302 readout}{c4}
\begin{scope}[on background layer]
\node[dgcluster2,fit=(c1)(c1l)(c2)(c2l)(c3)(c3l)(c4)(c4l)] (cc) {};\end{scope}
\node[dgctitle2,anchor=south west] at ([yshift=1.5mm]cc.north west) {Clinical source};
% cluster 3, simulation
\dgnodeg{dgtileg}{5.55}{0.55}{\glyphcpu}{QSP metastatic PDAC}{s1}
\dgnodeg{dgtileg}{7.75}{0.55}{\glyphdb}{100000 patient in silico}{s2}
\dgnodeg{dgtileg}{5.55}{-1.35}{\glyphai}{Digital twin proposals}{s3}
\dgnodeg{dgtileg}{7.75}{-1.35}{\glyphgear}{VVUQ test pipeline}{s4}
\begin{scope}[on background layer]
\node[dgcluster2,fit=(s1)(s1l)(s2)(s2l)(s3)(s3l)(s4)(s4l)] (cs) {};\end{scope}
\node[dgctitle2,anchor=south west] at ([yshift=1.5mm]cs.north west) {Simulation source};
% cluster 4, device
\dgnode{dgtile}{5.55}{-3.75}{\glyphrobot}{Eight-arm robotic stack}{d1}
\dgnode{dgtile}{7.75}{-3.75}{\glyphsignal}{Heartbeat broadcast bus}{d2}
\dgnode{dgtile}{6.65}{-5.65}{\glyphstop}{Force and no-fly limits}{d3}
\begin{scope}[on background layer]
\node[dgcluster,fit=(d1)(d1l)(d2)(d2l)(d3)(d3l)] (cd) {};\end{scope}
\node[dgctitle,anchor=south west] at ([yshift=1.5mm]cd.north west) {Device source};
% subject
\dgnodew{dgtiled}{10.60}{-2.55}{\glyphdoc}{IND v1.0, twelve modules}{ind}
\draw[trialburg,line width=0.6pt] (ind.center) circle (7.5mm);
% composed cluster
\dgnode{dgtilem}{13.85}{-0.55}{\glyphshield}{Physical AI Subpart J text}{x1}
\dgnode{dgtilem}{13.85}{-2.55}{\glyphuser}{Autonomy disclosure}{x2}
\dgnode{dgtilem}{13.85}{-4.55}{\glyphbank}{Verification cost ledger}{x3}
\begin{scope}[on background layer]
\node[dgcluster,fit=(x1)(x1l)(x2)(x2l)(x3)(x3l)] (cx) {};\end{scope}
\node[dgctitle,anchor=south west] at ([yshift=1.5mm]cx.north west) {Composed, no counterpart};
% edges, one per cluster, from a stated boundary waypoint
\draw[dgedgeb] (cr.east) -- node[above,font=\tiny\sffamily,fill=trialwhite,inner sep=1pt] {adapted} (ind.north west);
\draw[dgedgeb] (cd.east) -- node[below,font=\tiny\sffamily,fill=trialwhite,inner sep=1pt] {specified} (ind.south west);
\draw[dgedge] (cs.east) -- (ind.west);
\draw[dgedge] (cc.east) -- ++(0.55,0) -- (9.95,-4.35) -- (ind.south west);
\draw[dgedged] (ind.east) -- (x1.west);
\draw[dgedged] (ind.east) -- (x2.west);
\draw[dgedged] (ind.east) -- (x3.west);
\pnote{0}{-7.35}{Four clusters, one edge each, leaving from a stated waypoint on the cluster's
east boundary. The corridor between\\x = 8.95 and x = 10.15 carries no tile, so
the four inbound edges converge without crossing a label box.}
\end{appfig}
\vspace{-0.6cm}
\figcaption{Figure 6. Four source clusters feeding twelve IND modules, and the three\\
modules that had no prior-system counterpart and were composed instead.}
\end{appfloat}

\begin{apptable}
\begin{tabularx}{\textwidth}{>{\raggedright\arraybackslash}p{5.05cm} >{\raggedright\arraybackslash}p{1.55cm} >{\raggedright\arraybackslash}p{1.85cm} >{\raggedright\arraybackslash}X}
\headrow \thead{Module} & \thead{Position} & \thead{Characters} & \thead{What it establishes}\\
\midrule
Cover letter & 1 & 19,070 & Transmittal and sponsor identity \\
Table of contents & 2 & generated & ReGARDD content order \\
FDA forms 1571 and 1572 & 3 & 32,067 & The codified form set \\
Introduction & 4 & 82,102 & \S 312.23(a)(3)(i) introductory statement \\
General investigational plan & 5 & 76,989 & \S 312.23(a)(3)(iv) plan \\
Investigator's brochure & 6 & 47,368 & \S 312.23(a)(5) brochure \\
Proposed clinical research & 7 & 76,193 & \S 312.23(a)(6) protocols \\
Chemistry, manufacturing, control & 8 & 56,370 & \S 312.23(a)(7) CMC \\
Pharmacology and toxicology & 9 & 55,719 & \S 312.23(a)(8) nonclinical \\
Previous human experience & 10 & 48,045 & \S 312.23(a)(9) prior exposure \\
Additional information & 11 & 47,087 & \S 312.23(a)(10) \\
Relevant information and back matter & 12 & 53,647 & \S 312.23(a)(11) and references \\
\end{tabularx}
\end{apptable}
\vspace{-0.6cm}
\tabcap{Table 8. The twelve modules of the deposited IND in ReGARDD content order,\\their
character count at deposit, and the codified requirement each satisfies.}

\subsection{The clock}

The dossier was built in four stages over four days: a figure catalog, a draft,
a full version, and a senior-author final pass, deposited July 1, 2026 at
repository version v4.3.0 \cite{phase1ind}. The prior system's published
calendar for the same twelve modules is measured in months, and its stages carry
different names: regulatory drafting, nonclinical write-up, chemistry
compilation, internal quality control and legal review, then sponsor sign-off.
Figure 7 places both regimes on one axis with a single break glyph where the two
scales part, and Table 9 gives the day each module closed against the prior
system's name for the same work.

% ---------------------------------------------------------------------------
% FIGURE 7.  mermaid-type, gantt.  Two bands at one 0.35 cm row pitch and one
% 0.24 cm bar height, joined by a full-height break glyph, because 4 days and
% 390 days cannot share a linear scale legibly.
% ---------------------------------------------------------------------------
\begin{appfloat}[!tb]
\begin{appfig}[mermaid-type / gantt]
\ptitle{0}{0.95}{New system: four days, one author}
\axisx{0}{6.20}{0.35}
\foreach \x/\l in {0/{hour 0},1.55/{24},3.10/{48},4.65/{72},6.20/{96}}{
  \draw[trialslate,line width=0.3pt] (\x,0.35) -- (\x,0.45);
  \node[font=\tiny\sffamily,anchor=south] at (\x,0.47) {\l};}
\foreach \i/\lab/\x0/\x1 in {%
  0/{Cover letter and FDA 1571}/0/1.55, 1/{Introduction and plan}/0/1.55,
  2/{Investigator brochure}/1.55/3.10, 3/{Proposed clinical research}/1.55/3.10,
  4/{Chemistry manufacturing control}/1.55/3.10, 5/{Pharmacology and toxicology}/3.10/4.65,
  6/{Previous human experience}/3.10/4.65, 7/{Additional and relevant info}/3.10/4.65,
  8/{Figures, tables, back matter}/4.65/6.20}{
  \node[anchor=east,font=\tiny\sffamily] at (-0.15,{-0.35*\i}) {\lab};
  \fill[trialburgp] (\x0,{-0.35*\i-0.12}) rectangle (\x1,{-0.35*\i+0.12});
  \draw[trialburg,line width=0.3pt] (\x0,{-0.35*\i-0.12}) rectangle (\x1,{-0.35*\i+0.12});}
\node[anchor=east,font=\tiny\sffamily] at (-0.15,-3.15) {Defect pass and deposit};
\fill[trialburgl] (4.65,-3.27) rectangle (6.20,-3.03);
\draw[trialburg,line width=0.3pt] (4.65,-3.27) rectangle (6.20,-3.03);
% break glyph
\draw[trialchar,line width=0.9pt] (6.70,0.55) -- (6.90,-4.55);
\draw[trialchar,line width=0.9pt] (6.90,0.55) -- (7.10,-4.55);
% lower band
\ptitle{7.30}{0.95}{Prior system: the same work, in months}
\axisx{7.30}{13.50}{-4.05}
\foreach \x/\l in {7.30/{month 0},8.85/{3},10.40/{6},11.95/{9},13.50/{12}}{
  \draw[trialslate,line width=0.3pt] (\x,-4.05) -- (\x,-3.95);
  \node[font=\tiny\sffamily,anchor=south] at (\x,-3.93) {\l};}
\foreach \i/\lab/\x0/\x1 in {%
  0/{Regulatory drafting}/7.30/9.36, 1/{Nonclinical write up}/8.85/10.40,
  2/{CMC compilation}/9.62/10.91, 3/{Internal QC and legal}/10.40/11.43,
  4/{Sponsor sign off and deposit}/11.43/12.20}{
  \node[anchor=east,font=\tiny\sffamily] at (7.15,{-4.40-0.35*\i}) {\lab};
  \fill[trialmist] (\x0,{-4.40-0.35*\i-0.12}) rectangle (\x1,{-4.40-0.35*\i+0.12});
  \draw[trialslate,line width=0.3pt] (\x0,{-4.40-0.35*\i-0.12}) rectangle (\x1,{-4.40-0.35*\i+0.12});}
% ratio callout, in the empty right quadrant of the upper band
\node[mmgoal,text width=34mm] at (11.10,-1.55) {Ten modules in 96 hours\\against five stages in\\about 390 days};
% legend
\legkey{0}{-6.45}{trialburgp}{new system module}
\legkey{3.10}{-6.45}{trialburgl}{defect pass and deposit}
\legkey{6.60}{-6.45}{trialmist}{prior system stage}
\pnote{0}{-7.05}{Both bands use one 0.35 cm row pitch and one 0.24 cm bar height, so the two
regimes are visually commensurable\\even though their axes are not. The break
glyph is drawn once, at full canvas height, so it cannot be missed.}
\end{appfig}
\vspace{-0.6cm}
\figcaption{Figure 7. Twelve IND modules on one clock: the new system in hours
across four\\days, and the prior system's published calendar for the same work, in
months.}
\end{appfloat}

\begin{apptable}
\begin{tabularx}{\textwidth}{>{\raggedright\arraybackslash}p{3.9cm} >{\raggedright\arraybackslash}p{4.3cm} >{\raggedright\arraybackslash}p{1.5cm} >{\raggedright\arraybackslash}X}
\headrow \thead{Module} & \thead{Repository source read} & \thead{Window} & \thead{Prior-system counterpart}\\
\midrule
Introduction and plan & Adapted 21 CFR 312, ICH E6(R3) & Day 1 & Regulatory drafting, months \\
Investigator's brochure & Phase 1 protocol, RASolute 302 & Day 2 & Nonclinical write-up \\
Proposed clinical research & Phase 1 and Phase 2 protocols & Day 2 & Regulatory drafting \\
Chemistry, manufacturing, control & Public daraxonrasib record & Day 2 & CMC compilation \\
Pharmacology and toxicology & QSP and digital twin record & Day 3 & Nonclinical write-up \\
Previous human experience & RASolute 302, Dutch cohort & Day 3 & Internal quality control \\
Figures, tables, back matter & The 22-figure catalog & Day 4 & Sponsor sign-off \\
\end{tabularx}
\end{apptable}
\vspace{-0.6cm}
\tabcap{Table 9. Each IND module, the repository source the generation read,\\
the day it closed, and the prior system's name for the same work.}

\subsection{Complexity and volume}

The complexity feat is not the speed on its own. It is that the speed was
reached while the dossier stayed internally consistent across 594,657 characters
of regulatory prose, 52 distinct 21 CFR cross-references, 22 figures in a single
document sequence, 84 section-scoped tables, and 52 bibliography entries, with
every figure drawn in vector form and no raster anywhere. A reviewer opening the
IND at \S 312.23(a)(8) and following its cross-reference into the general
investigational plan finds the same dose levels, the same stopping rules and the
same device limits stated the same way, because all twelve modules were written
against one another rather than by separate hands. Figure 8 is the crosswalk
that makes that checkable, and Table 10 states the volume.

% ---------------------------------------------------------------------------
% FIGURE 8.  d2-type, sql tables.  Three record boxes at one 0.46 cm row height,
% two straight join edges leaving and entering at a row boundary.
% ---------------------------------------------------------------------------
\begin{appfloat}[!tb]
\begin{appfig}[d2-type / sql tables]
\def\rh{0.46}
\node[d2cellh,minimum width=4.6cm,minimum height=0.52cm,anchor=north west] at (0,0) {21 CFR 312.23 content item};
\foreach \i/\f in {1/{item id (PK)},2/{(a)(1) forms},3/{(a)(3)(i) introductory statement},
  4/{(a)(3)(iv) general plan},5/{(a)(5) investigator brochure},6/{(a)(6) protocols},
  7/{(a)(7) chemistry manufacturing},8/{(a)(8) pharmacology toxicology},
  9/{(a)(9) previous human experience},10/{(a)(10) and (a)(11) information}}{
  \node[d2celll,minimum width=4.6cm,minimum height=\rh cm,anchor=north west,align=left]
    at (0,{-0.52-\rh*(\i-1)}) {\quad \f};}
\node[d2cellh,minimum width=4.4cm,minimum height=0.52cm,anchor=north west] at (5.20,0) {Generated section file};
\foreach \i/\f in {1/{file id (PK)},2/{item id (FK)},3/{sec-NN-name.tex},
  4/{character count at deposit},5/{figures carried}}{
  \node[d2celll,minimum width=4.4cm,minimum height=\rh cm,anchor=north west,align=left]
    at (5.20,{-0.52-\rh*(\i-1)}) {\quad \f};}
\node[d2cellh,minimum width=4.4cm,minimum height=0.52cm,anchor=north west] at (10.20,0) {Repository deposit};
\foreach \i/\f in {1/{path id (PK)},2/{file id (FK)},3/{trial-ind/final-ind},
  4/{doi 10.5281/zenodo.21097442},5/{IND v1.0, repository v4.3.0}}{
  \node[d2celll,minimum width=4.4cm,minimum height=\rh cm,anchor=north west,align=left]
    at (10.20,{-0.52-\rh*(\i-1)}) {\quad \f};}
\draw[d2edgeb] (4.80,-1.21) -- node[above,font=\tiny\sffamily,fill=trialwhite,inner sep=1pt] {one to one} (5.20,-1.44);
\draw[d2edgeb] (9.60,-1.44) -- node[above,font=\tiny\sffamily,fill=trialwhite,inner sep=1pt] {one to one} (10.20,-1.44);
\node[d2soft,text width=44mm,anchor=north west] at (5.20,-3.85) {Twelve modules, 594657 characters};
\node[d2mid,text width=52mm,anchor=north west] at (10.20,-3.85) {22 figures, 84 section-scoped tables, 52 references};
\pnote{0}{-5.55}{The three record boxes share one 0.46 cm row height and one 0.52 cm header, so
each join edge leaves and enters\\at the vertical center of a row rather than at
an arbitrary point on a box edge.}
\end{appfig}
\vspace{-0.6cm}
\figcaption{Figure 8. Twelve codified IND content items joined to the section
file that\\satisfies each, and to the repository path where that file is
deposited.}
\end{appfloat}

\begin{apptable}
\begin{tabularx}{\textwidth}{>{\raggedright\arraybackslash}p{4.9cm} >{\raggedright\arraybackslash}p{2.9cm} >{\raggedright\arraybackslash}X}
\headrow \thead{Volume measure} & \thead{Deposited IND} & \thead{What it means for a reviewer}\\
\midrule
Section files & 12 & One per ReGARDD content position \\
Total characters across sections & 594,657 & The dossier's written weight \\
Figures, single document sequence & 22 & Every figure numbered once, end to end \\
Tables, section-scoped numbering & 84 & Table 3.1, Table 5.2, and so on \\
Distinct 21 CFR cross-references & 52 & Codified coverage that can be spot-checked \\
Bibliography entries & 52 & Every claim carries a source \\
Raster images & 0 & Every figure is vector and re-renderable \\
Authors & 1 & With model assistance, disclosed \\
Elapsed production time & 4 days & Against months in the prior system \\
\end{tabularx}
\end{apptable}
\vspace{-0.6cm}
\tabcap{Table 10. The volume one author reached in four days, measured on\\
nine axes, and what each measure means to a regulatory reviewer.}

Figure 9 states the sponsor's residual risk as logic rather than as a list. The
two clinical branches are AND gates, so a hold on clinical grounds requires both
a signal and a failure to respond to it, and the protocol's pause rule is
written to break that conjunction. The documentation branch is an OR gate, where
any one of four omissions suffices: insufficient chemistry information, a
brochure that omits a known risk, consent that lacks the autonomy disclosure, or
a sponsor report filed outside its window. That branch is exactly the class of
failure a system which can regenerate a complete, internally consistent dossier
in four days is built to eliminate, because the cost of a full re-issue after a
finding is four days rather than a quarter \cite{phase1ind, phase1guidance}.

% ---------------------------------------------------------------------------
% FIGURE 9.  graphviz-type, fault tree.  Four ranks at a uniform 2.15 cm
% separation, gate glyphs drawn rather than written, every basal event with
% exactly one parent so every minimal cut set is disjoint.
% ---------------------------------------------------------------------------
\begin{appfloat}[!tb]
\begin{appfig}[graphviz-type / fault tree]
\node[gvboxk,text width=44mm,line width=0.9pt] (top) at (7.50,0) {Clinical hold on the IND};
\umlgateor{7.50}{-2.15}
\umlgateand{2.20}{-4.30}
\umlgateand{7.50}{-4.30}
\umlgateor{12.80}{-4.30}
\node[gvnode,text width=24mm] (b1) at (0.95,-6.90) {Safety signal exceeds the protocol limit};
\node[gvnode,text width=24mm] (b2) at (3.45,-6.90) {Monitoring board does not pause in time};
\node[gvnode,text width=24mm] (b3) at (6.25,-6.90) {Device force cap breached in operation};
\node[gvnode,text width=24mm] (b4) at (8.75,-6.90) {Verification gate cleared on stale evidence};
\node[gvgray,text width=24mm] (b5) at (10.85,-6.90) {CMC information found insufficient};
\node[gvgray,text width=24mm] (b6) at (13.35,-6.90) {Brochure omits a known risk};
\node[gvgray,text width=24mm] (b7) at (10.85,-8.55) {Consent lacks the autonomy disclosure};
\node[gvgray,text width=24mm] (b8) at (13.35,-8.55) {Sponsor report outside its window};
\draw[gvedge] (top.south) -- (7.50,-1.89);
\draw[gvedge] (7.50,-2.41) -- (2.20,-4.04);
\draw[gvedge] (7.50,-2.41) -- (7.50,-4.04);
\draw[gvedge] (7.50,-2.41) -- (12.80,-4.04);
\draw[gvedge] (2.20,-4.56) -- (b1.north);
\draw[gvedge] (2.20,-4.56) -- (b2.north);
\draw[gvedge] (7.50,-4.56) -- (b3.north);
\draw[gvedge] (7.50,-4.56) -- (b4.north);
\draw[gvedge] (12.80,-4.42) -- (b5.north);
\draw[gvedge] (12.80,-4.42) -- (b6.north);
\draw[gvedge] (12.80,-4.42) -- (12.10,-5.85) -- (12.10,-8.55) -- (b7.east);
\draw[gvedge] (12.80,-4.42) -- (12.10,-5.85) -- (12.10,-8.55) -- (b8.west);
\node[gvboxm,text width=46mm] at (2.20,-8.55) {Two clinical branches are AND gates. The documentation branch is an OR gate, and is the branch the four-day re-issue removes.};
\pnote{0}{-9.85}{Each basal event has exactly one parent, so every minimal cut set is disjoint.
The only place an edge passes\\between two nodes is the 0.35 cm gap at x = 12.10,
and it is stated here for that reason.}
\end{appfig}
\vspace{-0.6cm}
\figcaption{Figure 9. The clinical hold as a top event over eight basal failures, with\\
the two AND gates that make a hold require a coincidence rather than a fault.}
\end{appfloat}

\subsection{What the volume did not cost}

Volume of this kind invites one obvious objection: that a dossier written at
this rate must be padded, and that 594,657 characters is a measure of verbosity
rather than of coverage. The objection is answerable, and the answer is in the
structure rather than in the prose. Every one of the twelve modules is written
to a codified content position, so a module cannot grow by adding material that
belongs somewhere else; the 52 distinct 21 CFR cross-references are checkable
one by one against the regulation; the 84 tables are section-scoped rather than
document-scoped, so a reviewer reading \S 312.23(a)(7) finds Table 6.1 and not
Table 43; and the 22 figures are numbered in a single document sequence, so a
cross-reference from the pharmacology module to a figure in the plan resolves
without ambiguity \cite{phase1ind}.

The second objection is that speed of this kind must come from copying. It
partly does, and the copying is disclosed: the proposed clinical research module
carries the Phase 1 protocol, the general investigational plan carries the Phase
2 forward path, and the adapted 21 CFR and ICH text supplies the regulatory
scaffolding, all of which are separately deposited and separately dated
\cite{onpremwhippl, phase2proto, natplatform2026}. What was not copied is the
part Figure 6 marks as composed: the Physical AI Subpart J text, the autonomy
disclosure, and the verification cost ledger. None of the three has a
prior-system counterpart to copy from, because no prior IND has been filed for a
device that generates the code it executes.

The third objection is the serious one: that a dossier assembled this fast
cannot have been checked. It was checked, twice, by two model manufacturers
independent of the one that produced it, at defined milestones and before
deposit, which is the review regime \S7 describes and measures
\cite{aiprgliomatc, rfarm27001v2}. That regime is not equivalent to a regulatory
review and does not claim to be. It is equivalent to the internal quality
control stage the prior system runs before sponsor sign-off, and it runs in
hours rather than in the sixty days that stage is typically budgeted.

\subsection{What a reviewer can check in an afternoon}

The practical test of a dossier is not its length but how quickly a reviewer can
falsify it. Four checks are available to anyone with the deposited files and no
special access. First, open Figure 8's crosswalk and pick any codified content
item; the named section file either satisfies it or it does not. Second, take
any of the 52 CFR cross-references and read the cited paragraph; the dossier
either states the requirement correctly or it does not. Third, take any number
that appears in more than one module, the day-30 safety window, the 3 N per-arm
force cap, the n = 18 ceiling, and check that every occurrence agrees. Fourth,
read the four-day production claim against the commit history of
\appfile{trial-ind}, which records when each module was written and pushed. Each
of the four is a falsification route, and each is open because the artifacts are
deposited rather than described.
